% Page 025

\seite{25}

Am Schlusse des Schuljahres wurden drei Kinder\ob
der Schule entlassen. Zwei Knaben.

\pend
\jahresueberschrift{Das Jahr 1897/98}
\pstart

Es wurden zu Ostern 1897 5 Schulerstlinge\ob
aufgenommen. Die Schülerzahl beträgt jetzt\ob
61.\ Sämtliche Kinder sind katholisch und voll\ohb
sinnig. Durch den starken Zuwachs werden die\ob
Schulräume zu eng, wodurch die Frage über einen\ob
Neubau immer mehr in den Vordergrund tritt.

\pend
\abschnitt{Pfarrerwechsel}
\pstart

Mit Beginn des Monats Mai wurde Herr Pfarrer\ob
Michael Dunckel, der auch zugleich Ortsschul\ohb
inspektor war, nach Hermeskeil versetzt. Die Pfarrei ließ den\ob
Herrn nur ungern scheiden, da er sich als echter\ob
Seelenhirt und wahrer Kinderfreund die Herzen\ob
aller im Sturme gewann. Sein Nachfolger\ob
wurde der bisherige Kaplan in St.\ Paulin-Trier\ob
Herr Ellges, der am 13.\ Mai 1897 feier\ohb
lich in sein Amt eingeführt wurde.

Mit Beginn des Sommersemesters trat ein Sohn d.\ob
hiesigen achtbaren Bürger ins Priesterseminar zu Trier\ob
ein, um sich dem erhabenen Berufe zu widmen.\ob
Seine Vorstudien hat er teils durch Privatunterricht,\ob
teils durch Besuch des Gymnasiums in Prüm erworben.\ob
Zwei junge Leute, die augenblicklich noch Privat\ohb
unterricht genießen, sind auch entschlossen, sich\ob
diesem Berufe zu widmen. Einer dieser entschloß sich noch
\pend
